\phantomsection
\section*{Appendix}
\addcontentsline{toc}{section}{Appendix}
\setcounter{section}{0}
\setcounter{subsection}{0}

\section{Author List}

\paragraph*{Core Contributors} Dongfang Li, Xiaodong Luo, Ruoyu Sun, Xuhui Chen, Linyuan Qiu, Jian Meng, Zhengxuan Lu, Yiting Wang, Yucheng Xie, Tao Guo, Tianxiang Fang, Jing Li, Sihang Chen, Shihao Hong, Chang Liu, Weihua Dai, Zirong Zeng, Ziwei Zhu, Zhuohan Wang, Zhengjun Yue, Igor Vasilyev, Min Liu, Weijian Sun,  Xin Chen 
\paragraph*{Contributors} Yingmeng Gao, Jinhua Zhou, Taolue Chen, Chenwei Wu, Dong Zhang, Wenlong Jin, Jinmin Xiang, Barkova Maria, Ushakov Anton, Xianfei Jin, Tian Ding, Zhihang Lin, Qian Chen, Linxin Yang, Mingzhe Yang, Bingwei Zhang, Hongzhang Yang, Fangxue Zhang, Shijun Qin, Jie Yu, Cuihua Hu, Tolstykh Vasiliy, Nosov Ivan, Abdullin Amir, Zhichen Zhou, Xin Zhang, Zhixiong Ning, Xutong Zhao, Junjie Huang, Jiajun Liu, Weiyan Kong, Zheng Zhang, Wenhan Luo, Lin Hu, Yangbo Guo, Li Zeng, Shihao Zeng 
\paragraph*{Project Leaders} Baotian Hu, Min Zhang, Haizhou Li, Zhiquan Luo




\section{Post-training Recipe Details} \label{app:post_training_recipes} This appendix summarizes the main recipe-level settings used in the post-training experiments, focusing on configurations that affect optimization, parallel execution, memory usage, and reproducibility.

\paragraph{CPT recipe.} Table~\ref{tab:cpt_training_recipe} reports the CPT recipe used for the Ascend-based DeepSeek-V4 continued-pretraining run, including the long-context setup, optimization configuration, parallelism strategy, and memory-saving mechanisms.



\paragraph{SFT recipe.}
Table~\ref{tab:sft_training_recipe} reports the main recipe-level configuration
used for supervised fine-tuning. In contrast to CPT, which optimizes
autoregressive language modeling over domain-adaptive corpora, SFT optimizes
instruction-response trajectories and masks the loss on prompt tokens so that
gradients are applied only to the target response. The reported items are
selected to match the capabilities discussed in Section~\ref{sec:sft}, including
instruction-following alignment, solver-oriented modeling, and
pipeline-optimized distributed execution on Ascend NPUs.



% \begin{table}[t]
% \centering
% \small
% \begin{tabular}{p{0.30\linewidth}p{0.58\linewidth}}
% \toprule
% Category & Configuration \\
% \midrule
% Training stage & Supervised fine-tuning \\
% Training objective & Instruction-response SFT with prompt-token loss masking \\
% Sequence length & [to be specified] \\
% Global batch size & [to be specified] \\
% Optimizer & [to be specified] \\
% Peak learning rate & [to be specified] \\
% Learning-rate schedule & [to be specified] \\
% Warmup rate & [to be specified] \\
% Gradient clipping & [to be specified] \\
% Precision & [to be specified] \\
% Parallelism & [to be specified] \\
% Activation checkpointing & [to be specified] \\
% Training sample scale & [to be specified] \\
% \bottomrule
% \end{tabular}
% \caption{Main recipe-level configuration for the SFT run. The table reports
% settings that affect optimization behavior, reproducibility, and the
% instruction-tuning experiments discussed in Section~\ref{sec:sft}.}
% \label{tab:sft_training_recipe}
% \end{table}

\paragraph{Agentic rollout configuration.}
The agentic rollout pipeline is evaluated in inference mode and does not update
model parameters in this work. Its configuration is therefore separated from
the CPT and SFT training recipes. If reported, the rollout configuration should
include only parameters that affect inference-time behavior and future
reinforcement-learning reproducibility, such as the number of sampled
formulation candidates, maximum self-review rounds, maximum code-repair rounds,
solver timeout, execution success criterion, and objective-matching tolerance.
These settings are rollout parameters rather than optimizer or training
hyperparameters.

% \begin{table}[t]
% ...
% \caption{Main inference-time rollout configuration for the OR agentic
% pipeline.}
% \label{tab:agent_rollout_config}
% \end{table}


% \% ============================================================================
% Overleaf-ready appendix fragment for:
%   SLAI2 Training Technical Report (Part I)
%
% Intended placement:
%   Insert after Appendix A.1 (Post-training Recipe Details). It will be
%   numbered A.2 automatically.
%
% Required packages (add to the main preamble if not already loaded):
%   \usepackage{amsmath,amssymb}
%   \usepackage{booktabs,tabularx,array}
%   \usepackage{graphicx,xcolor,listings,float}
%   \usepackage{enumitem}
%
% Disclosure policy:
%   This appendix intentionally contains no production-run identifiers,
%   corpus counts, acceptance rates, throughput measurements, internal paths,
%   prompt hashes, or historical experiment records. All numerical examples
%   are manually constructed illustrations and are not sampled from the
%   production corpus.
% ============================================================================

\definecolor{ORGray}{RGB}{245,247,249}

\lstdefinestyle{orcptappendix}{
  basicstyle=\ttfamily\footnotesize,
  breaklines=true,
  columns=fullflexible,
  keepspaces=true,
  frame=single,
  rulecolor=\color{black!25},
  backgroundcolor=\color{ORGray},
  showstringspaces=false,
  xleftmargin=0.5em,
  xrightmargin=0.5em
}

\section{Solver-Verified OR-CPT Data Synthesis: Engine Design and Illustrative Cases}
\label{app:or-cpt-engine}

Section 3.2 introduced the synthesized component of the operations-research
(OR) corpus and its bidirectional construction pipeline.  This appendix focuses
on the method: how a parameterized optimization generator is converted into a
self-contained CPT document, how mathematical invariants are preserved across
language-model transformations, and how deterministic checks prevent plausible
but incorrect samples from entering the training corpus.

The central design principle is that language-model fluency is not used as a
proxy for mathematical correctness.  Language models perform two semantic
transformations---from a verified optimization instance to a business problem,
and from that business problem back to a mathematical model and executable
program.  Generator contracts, static checks, and independent solver execution
surround these transformations and determine whether the resulting document is
admissible for CPT.

\begin{table}[H]
\centering
\small
\caption{Main recipe-level configuration for the Ascend-based CPT run under
the DeepSeek-V4-Flash long-context setting. The table reports settings
that affect optimization, parallel execution, memory behavior, and
NPU-specific training support. Dataset composition, checkpoint paths, logging
options, profiling settings, and environment-specific launch details are
omitted.}
\label{tab:cpt_training_recipe}
\begin{tabular}{p{0.36\linewidth}p{0.52\linewidth}}
\toprule
Configuration item & Setting \\
\midrule
\multicolumn{2}{c}{\textit{Training Setup}} \\
\midrule
Training stage & Full-parameter continued pre-training \\
Sequence length & 65,536 tokens \\
Global batch size & 128 sequences \\
MTP & Enabled \\
\midrule
\multicolumn{2}{c}{\textit{Optimization}} \\
\midrule
Optimizer & AdamW \\
Peak learning rate & $1.0\times10^{-6}$ \\
Learning-rate schedule & Cosine decay \\
Warmup rate & 20\% \\
Minimum learning-rate factor & 0.01 \\
Optimizer epsilon & $1.0\times10^{-6}$ \\
\midrule
\multicolumn{2}{c}{\textit{Parallelism}} \\
\midrule
Tensor parallel degree & 1 \\
Pipeline parallel degree & 1 \\
Context parallel degree & 16 \\
Expert parallel degree & 32 \\
FSDP-style sharding & Enabled with reshard-after-forward \\
\midrule
\multicolumn{2}{c}{\textit{Memory and NPU Support}} \\
\midrule
Activation checkpointing & Full \\
Optimizer memory strategy & Swap optimizer \\
Model-specific NPU kernels & Adapted for attention and MoE computation \\
\bottomrule
\end{tabular}
\end{table}
\begin{table}[t]
\centering
\small
\caption{Main recipe-level configuration for the SFT run. The table reports settings that affect optimization behavior, reproducibility, and the instruction-tuning experiments discussed in Section~\ref{sec:sft}.}
\begin{tabular}{p{0.36\linewidth}p{0.52\linewidth}}
\toprule
Configuration item & Setting \\
\midrule
\multicolumn{2}{c}{\textit{Training Setup}} \\
\midrule
Training stage & Supervised fine-tuning\\
Sequence length & 8,192 tokens \\
Global batch size & 128 sequences \\
MTP & Enabled \\
\midrule
\multicolumn{2}{c}{\textit{Optimization}} \\
\midrule
Optimizer & AdamW \\
Peak learning rate & $5.0\times10^{-6}$ \\
Minimum learning rate & $5.0\times10^{-8}$ \\
Learning-rate schedule & Cosine decay \\
Weight decay & $1.0\times10^{-2}$ \\
Gradient clipping & Maximum norm of 1.0 \\
Initial loss scale & 65,536 \\
\midrule
\multicolumn{2}{c}{\textit{Parallelism}} \\
\midrule
Tensor parallel degree & 1 \\
Pipeline parallel degree & 4 \\
Context parallel degree & 1 \\
Context parallel algorithm & Ulysses CP \\
Expert parallel degree & 32 \\
Sequence parallel & Enabled \\
\midrule
\multicolumn{2}{c}{\textit{Memory and Execution}} \\
\midrule
Activation checkpointing & Full \\
Optimizer memory strategy & Swap optimizer \\
MTP memory-efficient logits & Enabled \\

\bottomrule
\end{tabular}
\label{tab:sft_training_recipe}
\end{table}
\subsection{Role in the post-training workflow}
\label{app:orcpt-role}

The OR-CPT Engine is an upstream data-construction system.  It is separate from
the Ascend 910C training stack: the engine establishes mathematical fidelity,
semantic coverage, and provenance, while the training stack tokenizes, packs,
samples, and optimizes the model on the resulting corpus.  The engine exports
plain self-contained documents for autoregressive continued pre-training rather
than instruction--response conversations.

This separation also clarifies the relationship between CPT and SFT.  CPT
documents expose the model to OR terminology, formulation patterns, solver APIs,
and verification habits.  SFT subsequently teaches the model to elicit that
knowledge under task instructions and output contracts.  The same OR taxonomy
can guide both stages, but the data representations and optimization objectives
remain distinct.

\subsection{End-to-end engine architecture}
\label{app:orcpt-architecture}

The engine follows a bidirectional verification path.  It first samples a
structured optimization instance from a parameterized generator, solves the
source model independently, and filters infeasible, degenerate, numerically
unstable, or duplicated seeds.  The verified seed is then backtranslated into a
self-contained business problem, where a natural-language quality filter checks
numeric coverage, unit consistency, business grounding, and solution leakage.
Finally, the engine reconstructs the model from the generated problem, applies
static and contract-based validation, executes the reconstructed model, compares
its status and objective against the independent reference, and renders accepted
pairs into CPT documents with lineage metadata.

For a seed $s$, let $M_s$ denote the source model and $z_s^\star$ its reference
objective.  Backtranslation yields a natural-language problem $q_s$, and forward
modeling yields an executable reconstruction $\widehat{M}_s$ with objective
$\widehat{z}_s$.  The core admission condition can be summarized as

\begin{equation}
  \operatorname{Accept}(s)=
  \mathbb{I}\!\left[
    Q_{\mathrm{seed}}(M_s)
    \land Q_{\mathrm{NL}}(q_s)
    \land Q_{\mathrm{contract}}(\widehat{M}_s)
    \land Q_{\mathrm{exec}}(\widehat{M}_s)
    \land Q_{\mathrm{obj}}(z_s^\star,\widehat{z}_s)
  \right].
  \label{eq:orcpt-acceptance}
\end{equation}


The objective comparison uses absolute and relative tolerances.  It is a strong
executable consistency signal, but not a proof of symbolic equivalence: two
different formulations can share the same optimum on a particular numerical
instance. The static and contract-based checks therefore remain necessary even
when the objective agrees. Figure~\ref{fig:orcpt-trust-path} summarizes the
resulting end-to-end trust path from a generated seed to an accepted,
solver-verified CPT document.

\begin{figure}[H]
  \centering
  \fbox{%
    \begin{minipage}{0.94\linewidth}
      \centering\small
      \textbf{Generator}\ $\rightarrow$\
      \textbf{verified seed}\ $\rightarrow$\
      \textbf{business problem}\ $\rightarrow$\
      \textbf{reconstructed model}\ $\rightarrow$\
      \textbf{solver-verified CPT document}

      \vspace{0.5em}
      \textit{Profiles and contracts preserve mathematical structure; scenario
      controls vary the semantic wrapper; solver execution anchors correctness.}
    \end{minipage}}
  \caption{The trust path of a synthesized OR-CPT document.}
  \label{fig:orcpt-trust-path}
\end{figure}

\subsection{Mathematical seeds and domain contracts}
\label{app:orcpt-contracts}

\paragraph{Seed representation.}
A seed is a structured optimization instance, not merely a textual question.
It contains declared sets, coefficient tables, objective direction, variable
domains, bounds, constraint families, and executable or LP artifacts.  It also
carries semantic metadata such as task family, sub-family, modeling concepts,
difficulty axes, and formulation variant.  Deterministic random seeds support
regeneration when a generator must be audited or refined.

\paragraph{Generator profiles and family contracts.}
A generator profile describes the canonical mathematical signature of one
generator, including variable types, objective terms, core constraints, valid
indices, and forbidden reinterpretations.  A family contract describes
invariants shared across related generators, such as balance equations, linking
rules, conservation laws, and required variable families.  These layers are
grounded in the source facts of the seed---declared sets, coefficient tables,
bounds, units, and executable artifacts---so that the language model cannot
invent numbers or refer to nonexistent indices.  Seed-quality records further
capture feasibility, objective recomputation, solution activity, degeneracy, and
numerical stability, while lineage metadata records the generator, scenario
configuration, prompt revision, validation status, and rendering status outside
the training text.

% \begin{table}[H]
%   \centering
%   \caption{Contract layers surrounding a generated seed.}
%   \label{tab:orcpt-contract-layers}
%   \small
%   \begin{tabularx}{\linewidth}{@{}p{0.21\linewidth}p{0.31\linewidth}X@{}}
%     \toprule
%     Layer & Representative content & Primary protection \\
%     \midrule
%     Source facts
%       & Index sets, coefficient tables, bounds, units, LP/model artifact
%       & Prevent numeric invention and use of nonexistent indices. \\
%     Generator profile
%       & Variable types, objective terms, core constraints, forbidden changes
%       & Prevent reinterpretation as a neighboring but different OR problem. \\
%     Family contract
%       & Required variable families, balance/linking rules, valid index policy
%       & Prevent omission of structural constraints and objective components. \\
%     Seed-quality record
%       & Solver feasibility, objective recomputation, activity and degeneracy signals
%       & Prevent trivial, unstable, or educationally weak instances. \\
%     Lineage metadata
%       & Generator, scenario axes, prompt revision, evaluation and rendering status
%       & Make accepted documents traceable without placing internal metadata in the training text. \\
%     \bottomrule
%   \end{tabularx}
% \end{table}

\paragraph{Seed-quality filtering.}
Feasibility alone does not guarantee training value.  A model in which all
meaningful decisions are zero, every decision is fixed at a bound, or no
constraint participates in a genuine trade-off may be solvable but unhelpful.
The engine therefore checks solution activity, constraint activity, objective
recomputation, numerical precision, and duplicate signatures before a seed is
eligible for language generation.

\subsection{Prompt-controlled bidirectional synthesis}
\label{app:orcpt-prompts}

Backtranslation and forward modeling have different information boundaries and
are controlled by complementary fidelity requirements.  Backtranslation may
inspect the verified source facts, but it must preserve coefficients, bounds,
units, and indexed relations while hiding solver status, objective values,
variable solutions, LP artifacts, and code.  It may vary the industry setting,
stakeholder perspective, planning trigger, horizon, naming style, and units, but
it must not add or remove mathematical structure.  Forward modeling receives the
generated problem and the applicable structural contracts; it must reconstruct
the variables, objective sense, core constraints, index policy, explanation,
formulation, standalone Gurobi code, and self-check using only the declared data,
without hard-coding the reference answer or inferring missing coefficients from
the optimum.

% \begin{table}[H]
%   \centering
%   \caption{Principal controls in the two language-model transformations.}
%   \label{tab:orcpt-prompt-controls}
%   \small
%   \begin{tabularx}{\linewidth}{@{}p{0.22\linewidth}X X@{}}
%     \toprule
%     Control & Backtranslation & Forward modeling \\
%     \midrule
%     Source fidelity
%       & Preserve all available coefficients, bounds, units, and indexed relations.
%       & Use only declared data structures and reproduce each coefficient family exactly once. \\
%     Structural fidelity
%       & Express the canonical signature in business language without adding constraints.
%       & Enforce variable domains, objective sense, core constraints, and index policy. \\
%     Leakage control
%       & Hide solution, objective value, solver status, LP text, and code.
%       & Do not hard-code the optimum or infer missing coefficients from it. \\
%     Scenario control
%       & Apply industry, stakeholder, trigger, horizon, entity, and unit guidance.
%       & Preserve the business-to-model mapping stated in the problem. \\
%     Output contract
%       & Return a self-contained problem and structured data summary.
%       & Return explanation, formulation, standalone Gurobi code, and self-check. \\
%     \bottomrule
%   \end{tabularx}
% \end{table}

The production prompts are versioned implementation artifacts.  Their essential
logic can be represented without reproducing their full text:

\begin{lstlisting}[style=orcptappendix,caption={Abstract prompt interface used by the bidirectional synthesis stage.}]
BACKTRANSLATE(verified_seed, profile, scenario):
    preserve source facts and mathematical structure
    hide reference solution and implementation artifacts
    return a complete business problem

FORWARD_MODEL(business_problem, profile, family_contract):
    reconstruct explanation, formulation, and executable code
    use only declared indices and coefficients
    return a structured self-check; do not hard-code the optimum
\end{lstlisting}

The self-check is treated as a claim made by the generated response, not as the
validator.  Static inspection and independent execution remain authoritative
because prose may mention a required constraint that the code omits or indexes
incorrectly.

\subsection{Controlled scenarios and CPT rendering}
\label{app:orcpt-scenarios}

The scenario catalog separates mathematical diversity from linguistic
diversity.  A base OR structure can be expressed through different industry
lenses, decision triggers, organization types, planning frames, entity-naming
styles, and unit conventions.  These controls change the surface narrative---for
example, whether the problem is framed as a healthcare allocation, logistics
planning, energy scheduling, manufacturing review, outage response, budget
revision, or service-network decision---while keeping the source coefficient
tables, variable roles, objective terms, constraints, and mathematical periods
invariant.

% \begin{table}[H]
%   \centering
%   \caption{Semantic variation axes used by the scenario system.}
%   \label{tab:orcpt-scenario-axes}
%   \small
%   \begin{tabularx}{\linewidth}{@{}p{0.22\linewidth}p{0.32\linewidth}X@{}}
%     \toprule
%     Axis & Illustrative values & Controlled property \\
%     \midrule
%     Industry lens & healthcare, logistics, energy, manufacturing
%       & Domain vocabulary and plausible entities. \\
%     Decision trigger & demand revision, outage, budget review
%       & Concrete reason for the planning problem. \\
%     Organization & public agency, plant, service network
%       & Decision owner and stakeholder perspective. \\
%     Planning frame & daily plan, monthly review, board packet
%       & Narrative horizon without inventing mathematical periods. \\
%     Entity naming & indexed labels, role-based names
%       & Surface form of sets while retaining traceability. \\
%     Unit style & service units, hours, tons, currency
%       & Consistency among quantities, capacities, and objective terms. \\
%     \bottomrule
%   \end{tabularx}
% \end{table}

Negative guidance is equally important.  A transportation matrix must not be
turned into vehicle routing unless vehicles and sequencing exist.  A facility-
location task must retain its fixed opening decisions.  A lot-sizing task must
retain its period-by-period inventory balance.  These restrictions prevent a
fluent scenario from silently changing the model family.

After executable verification, the renderer converts the accepted pair into a
self-contained CPT document.  A full document can contain: (i) the business
problem and data; (ii) a business-to-model explanation; (iii) a mathematical
formulation; (iv) standalone solver code; (v) an interpreted solution; and
(vi) feasibility and objective-recalculation checks.  Style and section order
may vary, but mandatory mathematical and validation content does not.

\subsection{Illustrative case I: capacitated facility location}
\label{app:orcpt-case-cflp}

Suppose a regional service operator considers three candidate sites serving
three demand zones.  Opening a site incurs a fixed cost; service can be allocated
only from opened sites and cannot exceed site capacity.  The illustrative data
are shown in Table~\ref{tab:orcpt-example-data}.

\begin{table}[H]
  \centering
  \caption{Manually constructed facility-location example.}
  \label{tab:orcpt-example-data}
  \small
  \begin{tabular}{@{}lrrrrr@{}}
    \toprule
    Site & Fixed cost & Capacity & Cost to Z1 & Cost to Z2 & Cost to Z3 \\
    \midrule
    A & 120 & 100 & 4 & 7 & 6 \\
    B &  95 &  65 & 6 & 3 & 5 \\
    C &  80 &  60 & 8 & 5 & 2 \\
    \midrule
    Demand & -- & -- & 40 & 35 & 25 \\
    \bottomrule
  \end{tabular}
\end{table}

Let $y_i\in\{0,1\}$ indicate whether site $i$ is opened and let
$x_{ij}\geq0$ denote units served from site $i$ to zone $j$.  A faithful model
is

\begin{align}
  \min_{x,y}\quad
    & \sum_{i\in I} f_i y_i
      + \sum_{i\in I}\sum_{j\in J} c_{ij}x_{ij},
      \label{eq:example-objective}\\
  \text{s.t.}\quad
    & \sum_{i\in I}x_{ij}=d_j,
      && \forall j\in J,\\
    & \sum_{j\in J}x_{ij}\leq u_i y_i,
      && \forall i\in I,\\
    & y_i\in\{0,1\},\quad x_{ij}\geq0.
\end{align}

The optimal illustrative plan opens B and C.  Site B serves 30 units of Z1 and
all 35 units of Z2; site C serves the remaining 10 units of Z1 and all 25 units
of Z3.  Site B is exactly at its capacity of 65.  The objective is

\begin{equation}
  95+80+6(30)+8(10)+3(35)+2(25)=590.
  \label{eq:example-recalculation}
\end{equation}

This example contains a genuine fixed-charge trade-off.  Site C has the lowest
fixed cost and the cheapest service to Z3, while B is cheapest for Z2.  Neither
can serve total demand alone under the stated capacities.  The model must jointly
choose sites and allocate service rather than select the cheapest coefficient in
each demand column.

The same mathematics could be rendered as clinic siting, regional repair depots,
or computing-service hubs.  The entities and units may change, but the three
invariants---fixed activation, demand satisfaction, and capacity linking---must
remain present.  Forward reconstruction is accepted only if it recovers these
invariants and reproduces the reference objective within tolerance.

\subsection{Illustrative case II: executable but structurally wrong models}
\label{app:orcpt-case-failure}

This example also shows why execution success is insufficient.
Table~\ref{tab:orcpt-example-errors} contrasts the faithful formulation with
two executable but structurally incorrect reconstructions and identifies the
validation gates that detect each error.

\begin{table}[H]
  \centering
  \caption{Illustrative structural errors and the gates that detect them.}
  \label{tab:orcpt-example-errors}
  \small
  \begin{tabularx}{\linewidth}{@{}p{0.22\linewidth}p{0.27\linewidth}p{0.19\linewidth}X@{}}
    \toprule
    Reconstruction & Structural defect & Illustrative optimum & Detection \\
    \midrule
    Faithful model
      & Fixed costs and capacity linking retained
      & 590
      & Accepted after contract and objective checks. \\
    Fixed costs omitted
      & Sites can be opened without paying activation cost
      & 315
      & Contract reports a missing objective term; objective comparison fails. \\
    Capacity constraints omitted
      & Site B serves all demand despite capacity 65
      & 565
      & Contract reports a missing constraint; objective comparison fails. \\
    \bottomrule
  \end{tabularx}
\end{table}

All three programs can be syntactically valid, executable, and solved to an
\texttt{OPTIMAL} status.  That status certifies only the optimum of the program
that was actually constructed.  It does not certify faithfulness to the
intended business problem.  The example therefore motivates three distinct
checks: static presence of required terms, family-level structural contracts,
and independent objective comparison.

Repair follows the same rule.  When a problem statement lacks a coefficient or
bound, the safe action is regeneration, rejection, or recovery from an
authoritative source table.  A repair must not invent a value merely to make the
model bounded or to approach the reference objective.

\section{CPT--SFT--Deployment--Evaluation Provenance}
\label{app:provenance}

Post-training large-scale MoE models involves multiple stages—CPT, SFT, weight conversion, serving, and evaluation—each producing artifacts dependent on upstream outputs. We introduce a lightweight provenance system that links all stages via per-stage manifests and a local file-based registry, to avoid opacity between benchmark scores and training configurations, checkpoints, or deployed weights, thereby ensuring clear attribution of improvements, reproducibility, and auditability of evaluation claims. The system serves three functions: (i) it enforces a verifiable chain from training data to scores, tracing each result to its exact checkpoint; (ii) it enables root-cause analysis of regressions by tracing deployed models back through conversion and training to the upstream CPT run; (iii) it provides programmatic support for automated tracking, dashboard monitoring, and artifact management, without database dependencies or disruption to existing workflows.

\subsection{Design principles}

\begin{enumerate}
  \item \textbf{Model identity is decoupled from the service endpoint.}
  A served model name, IP address, and port identify the inference entry point,
  not the model version. The authoritative model identity is carried by the
  artifact manifest at the weight directory, together with the upstream training
  lineage recorded at each stage. This separation allows multiple model versions
  to be served under the same endpoint name while preserving unambiguous
  identification of which weights are in use.

  \item \textbf{Every stage records its ground truth at the point of execution.}
  Each stage writes the actual paths, key environment variables, and upstream
  relationships into a structured manifest immediately upon completion. This
  manifest, rather than any later reconstruction, serves as the authoritative
  record for that stage. The point-of-execution recording eliminates the
  fragility of post-hoc documentation and ensures that environmental details
  (e.g., CANN version, parallelism configuration, tokenizer path) are captured
  before they can be altered or forgotten.

  \item \textbf{Historical experiments can be incorporated post-hoc.}
  Registration tooling accepts existing configuration snapshots, checkpoint
  directories, and converted artifacts, allowing runs completed before the
  provenance system was deployed to be retroactively linked into the provenance
  graph. This ensures that valuable legacy experiments are not excluded from
  cross-run comparison and trend analysis.
\end{enumerate}

\subsection{Provenance chain and registry layout}

The end-to-end chain spans five stages: CPT training, SFT training, HF artifact
conversion, model deployment, and benchmark evaluation. Each stage produces a
manifest file at its output location and registers a copy in a centralized JSON
registry. The registry is organized by stage type (training runs, artifacts,
deployments, evaluation runs), with a lightweight index file maintaining
aggregate counts. In addition to the centralized entries, manifests are always
co-located with the artifacts they describe---the HF artifact manifest resides
in the weight directory, the deployment manifest resides in the serving log
directory, and the evaluation provenance resides in the benchmark summary
directory. This dual placement ensures that the provenance record travels with
the artifact even if the centralized registry is unavailable, and that any
downstream consumer of the artifact (e.g., a separate evaluation pipeline) can
discover its lineage without consulting the central registry.

\subsection{Stage-level linkage mechanism}

\paragraph{CPT $\rightarrow$ SFT.}
Both CPT and SFT runs produce a configuration snapshot and lineage manifest
through a shared registration interface. The parent--child relationship is
established by matching the SFT checkpoint load path against the CPT checkpoint
save path, recording \texttt{linked} status with \texttt{path\_match} confidence
on success, or \texttt{missing} otherwise. This link is essential for
CPT-to-SFT transfer analysis: without it, a reported transfer gain cannot be
traced to the specific CPT checkpoint that provided the initialization.

\paragraph{SFT $\rightarrow$ HF artifact.}
Upon weight format conversion to HuggingFace-compatible shards, the converter
writes an artifact manifest recording the artifact identifier, source training
run, checkpoint directory and iteration, HF output directory, and upstream CPT
lineage, with a copy registered centrally. This manifest is the authoritative
record connecting deployable weights to their training provenance.

\paragraph{HF artifact $\rightarrow$ model deployment.}
The serving launch procedure reads the artifact manifest from the weight
directory and generates a deployment manifest capturing the deployment
identifier, artifact and training run references, model path, served model name,
and inference endpoint URL, with copies in both the log directory and the
central registry. This artifact-to-endpoint binding enables the dashboard to
track which model version is served at each endpoint.

\paragraph{Model deployment $\rightarrow$ eval run.}
Evaluation accepts an optional deployment manifest reference via an environment
variable. When provided, the harness extends the benchmark summary with full
provenance fields spanning evaluation, deployment, artifact, SFT, and CPT.
Without it, evaluation proceeds normally but scores cannot be automatically
traced to specific weights or associated with the corresponding dashboard
experiment.

\subsection{Confidence levels}

Each provenance link carries a confidence annotation, summarized in
Table~\ref{tab:provenance_confidence}, that distinguishes
manifest-based exact linkage from path-based or heuristic matching.

\begin{table}[t]
\centering
\small
\caption{Provenance confidence levels assigned to each link in the chain.}
\label{tab:provenance_confidence}
\begin{tabular}{lp{0.65\linewidth}}
\toprule
Level & Condition \\
\midrule
\texttt{exact\_manifest}
  & The downstream stage directly reads an explicit manifest produced by the
  upstream stage (e.g., the deployment script reads the artifact manifest
  from the HF weight directory). \\
\texttt{path\_match}
  & The link is established by matching path fields across independently
  produced configuration records (e.g., the SFT checkpoint load path matches
  the CPT checkpoint save path). \\
\texttt{inferred\_from\_log}
  & Reserved for future use, where training or deployment logs may provide
  auxiliary evidence for a link. \\
\texttt{missing}
  & Insufficient information exists to establish the link; the downstream
  stage operates without a confirmed upstream connection. \\
\bottomrule
\end{tabular}
\end{table}

\subsection{Dashboard monitoring and retroactive registration}

The provenance registry directly feeds a monitoring dashboard that visualizes the
full experiment graph, each row corresponds to a training run, and
columns display the run identifier, job type, training iterations, key metrics
(loss, benchmark scores), and---critically---provenance status and confidence
annotations for each link in the chain. Runs with complete manifest coverage are
marked with \texttt{exact\_manifest} or \texttt{path\_match} confidence; legacy
runs that predate the provenance system are displayed with a \texttt{missing}
annotation but are otherwise fully functional in the dashboard.

The dashboard is refreshed by a build script that scans the training archive
directory and the centralized registry, aggregating manifest contents into a
single JSON consumed by the frontend. Newly completed runs appear automatically
on the next rebuild; legacy runs without manifests can be registered
retroactively by pointing the registration tooling at their existing
configuration files, checkpoint directories, or artifact outputs. This design
ensures that the dashboard reflects both actively tracked experiments and
backfilled historical data within a unified view, making it the operational
interface through which experiment provenance is queried, audited, and
cross-referenced against benchmark claims.

% ============================================================
\section{Training Schedule and Monitoring Design}
\label{sec:curriculum}

While the data refinement pipeline improves individual sample quality, the
macro-level organization of training data also requires careful design.
Our data-mixture analysis shows that overly concentrated OR-domain adaptation
can weaken general capabilities, motivating a training schedule that balances
OR, mathematics, code, and general-domain data. We therefore design a
predefined curriculum schedule together with a monitoring dashboard for
tracking training stability and capability retention. The schedule, stage
boundaries, and transition rules are specified before training, while the
dashboard serves as an observability tool for manual inspection.


% \subsection{Motivation: Examples of Potential Degradation}

% During development, we observed several failure patterns on an auxiliary code-generation benchmark that illustrate why a deliberate training schedule is necessary. These examples are representative of the fragility that can emerge when domain-specific data dominates the training distribution.

% \textbf{Case 1: Indentation inconsistency.} A simulation problem produced code with mismatched indentation:
% \begin{center}
% \begin{minipage}{0.2\textwidth}
% \begin{verbatim}
% B = [0] * N
% ps = 0
% for i in range(N):
% \end{verbatim}
% \end{minipage}
% \end{center}
% The variable \texttt{ps = 0} is incorrectly indented, causing an \texttt{IndentationError}. This basic syntax failure cannot be compensated for by domain knowledge.

% \textbf{Case 2: Spurious module name corruption.} For a graph problem, the model output \texttt{from collections import dequeorate} instead of \texttt{deque}. This import error passes superficial syntax checks but fails at runtime.

% \textbf{Case 3: Superficially plausible but incorrect algorithm.} On a minimum-swaps problem, the model applied a heuristic and returned an incorrect answer, indicating reliance on memorized patterns rather than robust reasoning.

% These observations motivated the design of a training schedule that explicitly balances OR data with general code and math data, and the inclusion of code-quality metrics in our monitoring dashboard.

\subsection{Static Curriculum Learning Design}

We organize the training data mixture into a pre-defined, three-stage schedule. The schedule uses a mixing vector \(\bm{\alpha}(t) = [\alpha_{\text{OR}}, \alpha_{\text{math}}, \alpha_{\text{code}}, \alpha_{\text{gen}}]\) that is varied across training steps \(t\) according to a fixed plan:

\begin{enumerate}
    \item \textbf{Foundation stage (first 30\% of steps):} Heavy emphasis on mathematical reasoning and clean general code (\(\alpha_{\text{math}}=0.35, \alpha_{\text{code}}=0.35, \alpha_{\text{gen}}=0.20, \alpha_{\text{OR}}=0.10\)). This phase reinforces precise token generation and algorithmic logic.

    \item \textbf{Domain ramp-up (30--70\%):} \(\alpha_{\text{OR}}\) is linearly increased to \(0.40\) while \(\alpha_{\text{code}}\) gradually reduces to \(0.15\). Within this stage, OR samples are introduced by difficulty format: first Data-in-Problem (DP), then Data-in-Table (DT), and finally Data-Problem-Separate (DPS). The transition to a more difficult format is gated by a fixed evaluation threshold on a held-out subset of the current format: the model must achieve more than \(85\%\) execution success before the next format is unlocked. This threshold is a hard rule, not an adaptive algorithm.

    \item \textbf{Contract hardening (70--100\%):} The mixture shifts back toward code-only and LP-writing tasks (\(\alpha_{\text{code}}\) raised to \(0.25\)) to reinforce strict output contracts, using validated samples from the data refinement pipeline.
\end{enumerate}

This schedule is a \textbf{default trajectory} whose parameters (proportions, thresholds, stage boundaries) were fixed based on preliminary runs and the data-mixture ablation in Section~\ref{subsec:sft_training_cleaning} (Figure~\ref{fig:ablation}). It is not modified dynamically during training. Curriculum-like scheduling strategies have been explored in other contexts~\citep{tzannetos2026curriculum,zhao2026rethinking,liang2025boosting}; our specific instantiation is one manually designed configuration.

\subsection{Monitoring Dashboard}

To observe training behavior and detect potential problems early, we implemented a monitoring setup that logs two groups of metrics at regular intervals (every 500 training steps). This dashboard is not a closed-loop controller; it is an observability layer for human developers.

\textbf{Training-health track.} Standard optimization signals: loss, gradient norm, NaN counts, and Model FLOPs Utilization (MFU). A spike in gradient norm or a NaN loss immediately signals a data or training issue requiring manual investigation.

\textbf{Capability-retention track.} A set of metrics that track both domain and general performance:
\begin{itemize}
    \item \textbf{OR capability:} validation perplexity on an OR-domain set (partitioned by format), and task-specific rewards (code\_reward, wl\_reward) on small proxy subsets of NL4OPT and OptiBench.
    \item \textbf{General code quality:} a \textit{syntax health score} computed on a held-out set of 200 pure-Python algorithmic problems. The score checks for non-ASCII tokens, indentation errors, import correctness, and basic static validity.
    \item \textbf{General language fluency:} perplexity on a fixed general-domain corpus (Wikipedia, books) to detect catastrophic forgetting.
\end{itemize}

These metrics are logged and visualized in WandB dashboards. During development, they were used for manual inspection: for example, a persistent decline in syntax health could prompt the developer to pause OR-centric training, inspect recent data, or adjust the mixture. The specific thresholds and the resulting decisions were made by human engineers, not by an automated state machine. The dashboard design is inspired by systematic analyses of training dynamics~\citep{qi2026evolm}, and its instrumentation code is included in the released artifacts.




% % ---------- Required packages ----------
% % \usepackage[most]{tcolorbox}
% % \usepackage{booktabs}
% % \usepackage{listings}
% % \usepackage{xcolor}

% % ---------- Color palette centered on #B34A78 ----------
% \newpage
% \section{Representative Original-Checkpoint Error Cases}
% \label{app:original_checkpoint_error_cases}

% Section~\ref{subsec:original_checkpoint_diagnosis} summarizes the aggregate
% failure distribution of the original DeepSeek-V4-Flash checkpoint.  This
% appendix records the concrete examples used in that diagnosis.  The cases are
% traced back to the corresponding evaluation artifacts and are grouped by the
% six failure types discussed in the Inference-Setting Shift analysis.

% % ---------- Color palette centered on #B34A78 ----------
% % \newpage
% \definecolor{casePrimary}{HTML}{B34A78}
% \definecolor{caseDark}{HTML}{7A2E50}
% \definecolor{caseLight}{HTML}{F8EDF2}
% \definecolor{caseMid}{HTML}{D99AB7}
% \definecolor{caseText}{HTML}{2B2227}
% \definecolor{caseMuted}{HTML}{7B6A72}
% \definecolor{errRed}{HTML}{B21F3A}

% \lstdefinestyle{pythoncase}{
%   language=Python,
%   basicstyle=\ttfamily\scriptsize,
%   keywordstyle=\color{caseDark}\bfseries,
%   commentstyle=\color{caseMuted},
%   stringstyle=\color{casePrimary},
%   breaklines=true,
%   columns=fullflexible,
%   keepspaces=true,
%   showstringspaces=false,
%   escapeinside={(*@}{@*)}
% }

% \begin{tcolorbox}[
%   colback=caseLight,
%   colframe=caseDark,
%   coltitle=white,
%   colbacktitle=casePrimary,
%   title={Case Study: Canonical LP-Graph Mismatch under ORGEval},
%   fonttitle=\bfseries,
%   sharp corners,
%   boxrule=0.9pt,
%   left=1.2mm,
%   right=1.2mm,
%   top=1mm,
%   bottom=1mm
% ]

% \begin{tabular}{p{0.22\linewidth}p{0.68\linewidth}}
% \toprule
% Benchmark & B4O ORGEval \\
% Sample ID & \texttt{BENCH4OPT\_0}; related: \texttt{BENCH4OPT\_1}, \texttt{BENCH4OPT\_12} \\
% Evaluation & \texttt{code\_reward=1.0}, \texttt{wl\_reward=0.0}; \texttt{var\_num\_check=True}, \texttt{cons\_num\_check=False} \\
% Failure Type & Executable LP with non-reference constraint family \\
% \bottomrule
% \end{tabular}

% \vspace{0.8em}
% \textbf{Original Problem}

% \small
% The problem is a pharmaceutical blending task.  The model must choose
% ingredient quantities to minimize cost while respecting ingredient
% availability, medication dosage requirements, and active-ingredient purity
% bounds.  ORGEval does not only check whether the generated LP can be written;
% it compares the variable family, constraint family, and their graph structure
% against the reference LP.

% \vspace{0.8em}
% \textbf{Model Output}

% \begin{lstlisting}[style=pythoncase]
% x = model.addVars(I, lb=0.0, ub=GRB.INFINITY, name='x')
% model.setObjective(gp.quicksum(costs[i] * x[i] for i in range(I)), GRB.MINIMIZE)

% for i in range(I):
%     model.addConstr(x[i] <= availability[i], name=f'availability_{i}')

% # The model assumes dosage bounds apply to the total medication amount.
% for m in range(M):
%     total_med = gp.quicksum(requirements[m][i] * x[i] for i in range(I))
%     model.addConstr(total_med >= min_dosage[m], name=f'min_dosage_{m}')
%     model.addConstr(total_med <= max_dosage[m], name=f'max_dosage_{m}')

% total_ingredients = gp.quicksum(x[i] for i in range(I))
% for i in range(I):
%     model.addConstr(x[i] >= min_purity[i] * total_ingredients,
%                     name=f'min_purity_{i}')
%     model.addConstr(x[i] <= max_purity[i] * total_ingredients,
%                     name=f'max_purity_{i}')
% \end{lstlisting}

% \vspace{0.4em}
% \textbf{Error Annotation.}
% The program is executable and writes an LP file, but the generated constraint
% family does not match the reference formulation.  In this family of cases,
% \texttt{BENCH4OPT\_1} first fails variable-count matching, and
% \texttt{BENCH4OPT\_12} shows the next level of difficulty: even when variable
% and constraint counts match, the WL graph can still be wrong.  These examples
% show why structural ORGEval failures motivate CPT for domain templates and SFT
% for benchmark-facing canonical formulations.

% \end{tcolorbox}

% \begin{tcolorbox}[
%   colback=caseLight,
%   colframe=caseDark,
%   coltitle=white,
%   colbacktitle=casePrimary,
%   title={Case Study: Output-Contract and Schema Failures},
%   fonttitle=\bfseries,
%   sharp corners,
%   boxrule=0.9pt,
%   left=1.2mm,
%   right=1.2mm,
%   top=1mm,
%   bottom=1mm
% ]

% \begin{tabular}{p{0.22\linewidth}p{0.68\linewidth}}
% \toprule
% Benchmark & B4O ORGEval \\
% Sample ID & \texttt{BENCH4OPT\_4}; \texttt{BENCH4OPT\_176} \\
% Evaluation & \texttt{BENCH4OPT\_4}: syntax failure from markdown fence; \texttt{BENCH4OPT\_176}: \texttt{TypeError} from list indexing with a string shift label \\
% Failure Type & Code-only protocol violation and JSON/list schema mismatch \\
% \bottomrule
% \end{tabular}

% \vspace{0.8em}
% \textbf{Original Problems}

% \small
% \texttt{BENCH4OPT\_4} is a capital-budgeting problem: select investment
% projects to maximize NPV under budget, risk, and dependency constraints.
% \texttt{BENCH4OPT\_176} is a hospital nurse-scheduling problem: assign nurses
% to shifts while satisfying shift coverage, skill requirements, and maximum
% working hours.

% \vspace{0.8em}
% \textbf{Model Output}

% \begin{lstlisting}[style=pythoncase]
% # BENCH4OPT_4: code-only target, but the generated response starts with:
% (*@\textbf{\textcolor{errRed}{```python}}@*)
% import json
% import gurobipy as gp
% from gurobipy import GRB

% # BENCH4OPT_176: shifts are string labels, while shift_requirements is a list.
% for j in shifts:
%     model.addConstr(
%         gp.quicksum(x[i, j] for i in nurses) >=
%         (*@\textbf{\textcolor{errRed}{\texttt{shift\_requirements[j]}}}@*),
%         name=f"shift_coverage_{j}"
%     )
% \end{lstlisting}

% \vspace{0.4em}
% \textbf{Error Annotation.}
% The first error blocks evaluation before the model is even built: the executor
% sees the markdown fence token as Python syntax.  The second error reaches
% model construction but fails when a string shift name is used as a list index.
% Both examples primarily motivate SFT: the model already has useful OR
% knowledge, but Pass@1 behavior must obey code-only contracts and schema-safe
% data-access templates.

% \end{tcolorbox}

% \begin{tcolorbox}[
%   colback=caseLight,
%   colframe=caseDark,
%   coltitle=white,
%   colbacktitle=casePrimary,
%   title={Case Study: Discrete Semantics Error in NL4OPT},
%   fonttitle=\bfseries,
%   sharp corners,
%   boxrule=0.9pt,
%   left=1.2mm,
%   right=1.2mm,
%   top=1mm,
%   bottom=1mm
% ]

% \begin{tabular}{p{0.22\linewidth}p{0.68\linewidth}}
% \toprule
% Benchmark & NL4OPT \\
% Sample ID & \texttt{NL4OPT\_15} \\
% Evaluation & Ground truth: 28250; Prediction: 28125 \\
% Failure Type & Continuous relaxation of integer workforce variables \\
% \bottomrule
% \end{tabular}

% \vspace{0.8em}
% \textbf{Original Problem}

% \small
% A store employs senior citizens who earn \$500 per week and young adults who earn \$750 per week.
% The store must keep the weekly wage bill below \$30000.
% On any day, the store requires at least 50 workers, of whom at least 10 must be young adults.
% To ensure the store runs smoothly, the number of young adults should be at least a third the number of senior citizens.
% Formulate a LP to minimize the wage bill.

% \vspace{0.8em}
% \textbf{Model Output}

% \begin{lstlisting}[style=pythoncase]
% import gurobipy as gp
% from gurobipy import GRB

% m = gp.Model("MinimizeWageBill")

% x = m.addVar(vtype=(*@\textbf{\textcolor{errRed}{GRB.CONTINUOUS}}@*), name="seniors")
% y = m.addVar(vtype=(*@\textbf{\textcolor{errRed}{GRB.CONTINUOUS}}@*), name="young_adults")

% m.setObjective(500 * x + 750 * y, sense=GRB.MINIMIZE)

% m.addConstr(x + y >= 50, "total_workers")
% m.addConstr(y >= 10, "min_young")
% m.addConstr(3 * y >= x, "young_senior_ratio")
% m.addConstr(500 * x + 750 * y <= 30000, "wage_budget")

% m.optimize()

% if m.status == GRB.OPTIMAL:
%     print(f"Optimal solution: x = {x.X:.2f}, y = {y.X:.2f}")
%     print(f"Minimum wage bill = ${m.ObjVal:.2f}")
% else:
%     print("No optimal solution found.")
% \end{lstlisting}

% \vspace{0.4em}
% \textbf{Error Annotation.}
% The problem asks for numbers of workers, but the model uses continuous variables.
% This relaxes the integer workforce decision and yields a fractional optimum.

% \end{tcolorbox}

% \begin{tcolorbox}[
%   colback=caseLight,
%   colframe=caseDark,
%   coltitle=white,
%   colbacktitle=casePrimary,
%   title={Case Study: Loss/Yield-Driven Conservation Error},
%   fonttitle=\bfseries,
%   sharp corners,
%   boxrule=0.9pt,
%   left=1.2mm,
%   right=1.2mm,
%   top=1mm,
%   bottom=1mm
% ]

% \begin{tabular}{p{0.22\linewidth}p{0.68\linewidth}}
% \toprule
% Benchmark & B4O Feasible \\
% Sample ID & \texttt{BENCH4OPT\_2} \\
% Evaluation & Pass@1 prediction: 265.5768398591736; ground truth: 531.1536797183472. 5-shot and Pass@16 recover the reference value. \\
% Failure Type & Evaporation loss mixed into both cost and flow conservation \\
% \bottomrule
% \end{tabular}

% \vspace{0.8em}
% \textbf{Original Problem}

% \small
% The task is a water-distribution network-flow problem.  Each edge has a
% capacity, transportation cost, and evaporation rate.  The model must satisfy
% node demand while minimizing the total cost induced by flow and evaporation.

% \vspace{0.8em}
% \textbf{Model Output}

% \begin{lstlisting}[style=pythoncase]
% flow = {}
% for i, (source, target, capacity, cost, evaporation_rate) in enumerate(edges):
%     flow[(source, target)] = model.addVar(lb=0.0, ub=capacity,
%                                          name=f"flow_{source}_{target}")

% obj = gp.LinExpr()
% for (source, target), var in flow.items():
%     ...
%     obj += (*@\textbf{\textcolor{errRed}{\texttt{(cost + evaporation\_rate) * var}}}@*)
% model.setObjective(obj, GRB.MINIMIZE)

% for node_idx, node in enumerate(nodes):
%     inflow = gp.LinExpr()
%     outflow = gp.LinExpr()
%     evaporation_loss = gp.LinExpr()
%     ...
%     model.addConstr(
%         inflow - outflow - (*@\textbf{\textcolor{errRed}{\texttt{evaporation\_loss}}}@*) == demand[node_idx],
%         name=f"flow_conservation_{node}"
%     )
% \end{lstlisting}

% \vspace{0.4em}
% \textbf{Error Annotation.}
% The code is syntactically valid and reaches an optimal solution, but it
% implements the wrong physical system.  Evaporation is treated as an additive
% unit cost and also deducted from conservation, blurring whether loss affects
% delivered quantity, incurred cost, or both.  This is a domain-pattern failure:
% network flow with loss, yield, shrinkage, and evaporation needs stable
% conservation templates rather than generic LP scaffolding.

% \end{tcolorbox}

% \begin{tcolorbox}[
%   colback=caseLight,
%   colframe=caseDark,
%   coltitle=white,
%   colbacktitle=casePrimary,
%   title={Case Study: Ratio, Unit, and Total-Quantity Coupling Error},
%   fonttitle=\bfseries,
%   sharp corners,
%   boxrule=0.9pt,
%   left=1.2mm,
%   right=1.2mm,
%   top=1mm,
%   bottom=1mm
% ]

% \begin{tabular}{p{0.22\linewidth}p{0.68\linewidth}}
% \toprule
% Benchmark & OptiBench \\
% Sample ID & \texttt{OPTIBENCH\_6} \\
% Evaluation & Pass@1 prediction: 10000.0; 5-shot prediction: 28000.0; Pass@16 prediction: 10000.0; ground truth: 18000.0 \\
% Failure Type & Investment-return units and total savings are coupled incorrectly \\
% \bottomrule
% \end{tabular}

% \vspace{0.8em}
% \textbf{Original Problem}

% \small
% A logistics company allocates 50 trucks among three goods.  Each good has an
% initial per-truck delivery cost.  Investment in route-optimization software can
% reduce per-truck cost, and total investment must not exceed the total savings
% generated by the investment.  The objective is to minimize total delivery cost
% while satisfying minimum truck counts for Goods X and Y.

% \vspace{0.8em}
% \textbf{Model Output}

% \begin{lstlisting}[style=pythoncase]
% x = m.addVars(goods, vtype=GRB.INTEGER, name="trucks")
% I = m.addVars(goods, vtype=GRB.CONTINUOUS, name="investment")

% # cost per truck = c_i - (r_i/100)*I_i
% obj = gp.QuadExpr()
% for i in goods:
%     obj += (c[i] - (r[i]/100.0) * I[i]) * x[i]
% m.setObjective(obj, GRB.MINIMIZE)

% m.addConstr(gp.quicksum(I[i] for i in goods) <= B, name="budget")

% # Savings constraint: total investment <= total cost savings
% savings_lhs = gp.quicksum(I[i] for i in goods)
% savings_rhs = gp.quicksum(
%     (*@\textbf{\textcolor{errRed}{(r[i]/100.0) * I[i] * x[i]}}@*) for i in goods
% )
% m.addConstr(savings_lhs <= savings_rhs, name="savings")
% m.params.NonConvex = 2
% \end{lstlisting}

% \vspace{0.4em}
% \textbf{Error Annotation.}
% The model recognizes that investment and truck allocation interact, but it does
% not correctly reconcile dollars, hundreds-of-dollars, per-truck reduction, and
% total savings.  In 5-shot, the error changes shape: the model moves toward a
% linear-looking savings template but can drop the scale variable.  This case
% therefore exposes a deeper semantic issue than API usage: ratio, ROI,
% average-cost, and efficiency statements require explicit tracking of numerator,
% denominator, unit, and total quantity.

% \end{tcolorbox}

% \begin{tcolorbox}[
%   colback=caseLight,
%   colframe=caseDark,
%   coltitle=white,
%   colbacktitle=casePrimary,
%   title={Case Study: Solver-Compatible Nonlinear Reformulation Failure},
%   fonttitle=\bfseries,
%   sharp corners,
%   boxrule=0.9pt,
%   left=1.2mm,
%   right=1.2mm,
%   top=1mm,
%   bottom=1mm
% ]

% \begin{tabular}{p{0.22\linewidth}p{0.68\linewidth}}
% \toprule
% Benchmark & OptiBench \\
% Sample ID & \texttt{OPTIBENCH\_13} \\
% Evaluation & Pass@1: optimization failed with ``Objective must be linear or quadratic''; 5-shot prediction: 0.125; Pass@16 prediction: 0.125; ground truth: 0.35355235026576626 \\
% Failure Type & Fourth-order objective passed directly to Gurobi \\
% \bottomrule
% \end{tabular}

% \vspace{0.8em}
% \textbf{Original Problem}

% \small
% A person is trapped near a swamp boundary described by \(y \ge x^2\) with
% \(x \ge 0\).  The task is to find the closest feasible rescue point to
% \((3/4,0)\), equivalently the shortest rope length to a point on or above the
% parabolic boundary.

% \vspace{0.8em}
% \textbf{Model Output}

% \begin{lstlisting}[style=pythoncase]
% m = gp.Model("SwampRescue")
% x = m.addVar(lb=0.0, vtype=GRB.CONTINUOUS, name="x")

% # Objective: minimize (x - 0.75)^2 + (x^2)^2
% obj = (x - 0.75)*(x - 0.75) + (*@\textbf{\textcolor{errRed}{x*x*x*x}}@*)
% m.setObjective(obj, GRB.MINIMIZE)

% m.optimize()
% \end{lstlisting}

% \vspace{0.4em}
% \textbf{Error Annotation.}
% The mathematical expression is meaningful, but it is not a linear or quadratic
% objective.  Later inference settings avoid the initial build error yet still
% return the wrong distance, so the missing capability is not merely a Gurobi API
% habit.  The model needs solver-aware reformulation knowledge: when to introduce
% auxiliary variables, when nonconvex quadratic or general constraints are
% legitimate, and when analytic simplification is required before coding the
% optimization model.

% \end{tcolorbox}


% ---------- Required packages ----------
% Put these in the preamble:
% \usepackage[most]{tcolorbox}
% \usepackage{booktabs}
% \usepackage{tabularx}
% \usepackage{listings}
% \usepackage{xcolor}

\newpage
\section{Representative Original-Checkpoint Error Cases}
\label{app:original_checkpoint_error_cases}

\raggedbottom

Section~\ref{subsec:original_checkpoint_diagnosis} summarizes the aggregate
failure distribution of the original DeepSeek-V4-Flash checkpoint. This
appendix records the concrete examples used in that diagnosis. The cases are
traced back to the corresponding evaluation artifacts and are grouped by the
six failure types discussed in the Inference-Setting Shift analysis.

% ---------- Color palette centered on #B34A78 ----------
\definecolor{casePrimary}{HTML}{B34A78}
\definecolor{caseDark}{HTML}{7A2E50}
\definecolor{caseLight}{HTML}{F8EDF2}
\definecolor{caseMid}{HTML}{D99AB7}
\definecolor{caseText}{HTML}{2B2227}
\definecolor{caseMuted}{HTML}{7B6A72}
\definecolor{errRed}{HTML}{B21F3A}

\lstdefinestyle{pythoncase}{
  language=Python,
  basicstyle=\ttfamily\scriptsize,
  keywordstyle=\color{caseDark}\bfseries,
  commentstyle=\color{caseMuted},
  stringstyle=\color{casePrimary},
  breaklines=true,
  columns=fullflexible,
  keepspaces=true,
  showstringspaces=false,
  aboveskip=0.4em,
  belowskip=0.2em,
  escapeinside={(*@}{@*)}
}

\newtcolorbox{casebox}[1]{
  enhanced,
  colback=caseLight,
  colframe=caseDark,
  coltitle=white,
  colbacktitle=casePrimary,
  title={#1},
  fonttitle=\bfseries,
  sharp corners,
  boxrule=0.9pt,
  boxsep=1mm,
  left=1mm,
  right=1mm,
  top=0.8mm,
  bottom=0.8mm,
  before skip=0pt,
  after skip=0pt
}

\newenvironment{casepage}[1]{%
  \clearpage
  \begin{casebox}{#1}
}{%
  \end{casebox}
}

\newenvironment{casepagestart}[1]{%
  \begin{casebox}{#1}
}{%
  \end{casebox}
}

\vspace{1em}
\begin{casepagestart}{Case Study: Canonical LP-Graph Mismatch under ORGEval}

\begin{tabularx}{\linewidth}{@{}p{0.22\linewidth}X@{}}
\toprule
Benchmark & B4O ORGEval \\
Sample ID & \texttt{BENCH4OPT\_0}; related: \texttt{BENCH4OPT\_1}, \texttt{BENCH4OPT\_12} \\
Evaluation & \texttt{code\_reward=1.0}, \texttt{wl\_reward=0.0}; \texttt{var\_num\_check=True}, \texttt{cons\_num\_check=False} \\
Failure Type & Executable LP with non-reference constraint family \\
\bottomrule
\end{tabularx}

\vspace{0.5em}
\textbf{Original Problem}

{\small
The problem is a pharmaceutical blending task. The model must choose
ingredient quantities to minimize cost while respecting ingredient
availability, medication dosage requirements, and active-ingredient purity
bounds. ORGEval does not only check whether the generated LP can be written;
it compares the variable family, constraint family, and their graph structure
against the reference LP.
\par}

\vspace{0.5em}
\textbf{Model Output}

\begin{lstlisting}[style=pythoncase]
x = model.addVars(I, lb=0.0, ub=GRB.INFINITY, name='x')
model.setObjective(gp.quicksum(costs[i] * x[i] for i in range(I)), GRB.MINIMIZE)

for i in range(I):
    model.addConstr(x[i] <= availability[i], name=f'availability_{i}')

# The model assumes dosage bounds apply to the total medication amount.
for m in range(M):
    total_med = gp.quicksum(requirements[m][i] * x[i] for i in range(I))
    model.addConstr(total_med >= min_dosage[m], name=f'min_dosage_{m}')
    model.addConstr(total_med <= max_dosage[m], name=f'max_dosage_{m}')

total_ingredients = gp.quicksum(x[i] for i in range(I))
for i in range(I):
    model.addConstr(x[i] >= min_purity[i] * total_ingredients,
                    name=f'min_purity_{i}')
    model.addConstr(x[i] <= max_purity[i] * total_ingredients,
                    name=f'max_purity_{i}')
\end{lstlisting}

\vspace{0.3em}
\textbf{Error Annotation.}
The program is executable and writes an LP file, but the generated constraint
family does not match the reference formulation. In this family of cases,
\texttt{BENCH4OPT\_1} first fails variable-count matching, and
\texttt{BENCH4OPT\_12} shows the next level of difficulty: even when variable
and constraint counts match, the WL graph can still be wrong. These examples
show why structural ORGEval failures motivate CPT for domain templates and SFT
for benchmark-facing canonical formulations.

\end{casepagestart}

\begin{casepage}{Case Study: Output-Contract and Schema Failures}

\begin{tabularx}{\linewidth}{@{}p{0.22\linewidth}X@{}}
\toprule
Benchmark & B4O ORGEval \\
Sample ID & \texttt{BENCH4OPT\_4}; \texttt{BENCH4OPT\_176} \\
Evaluation & \texttt{BENCH4OPT\_4}: syntax failure from markdown fence; \texttt{BENCH4OPT\_176}: \texttt{TypeError} from list indexing with a string shift label \\
Failure Type & Code-only protocol violation and JSON/list schema mismatch \\
\bottomrule
\end{tabularx}

\vspace{0.5em}
\textbf{Original Problems}

{\small
\texttt{BENCH4OPT\_4} is a capital-budgeting problem: select investment
projects to maximize NPV under budget, risk, and dependency constraints.
\texttt{BENCH4OPT\_176} is a hospital nurse-scheduling problem: assign nurses
to shifts while satisfying shift coverage, skill requirements, and maximum
working hours.
\par}

\vspace{0.5em}
\textbf{Model Output}

\begin{lstlisting}[style=pythoncase]
# BENCH4OPT_4: code-only target, but the generated response starts with:
(*@\textbf{\textcolor{errRed}{```python}}@*)
import json
import gurobipy as gp
from gurobipy import GRB

# BENCH4OPT_176: shifts are string labels, while shift_requirements is a list.
for j in shifts:
    model.addConstr(
        gp.quicksum(x[i, j] for i in nurses) >=
        (*@\textbf{\textcolor{errRed}{\texttt{shift\_requirements[j]}}}@*),
        name=f"shift_coverage_{j}"
    )
\end{lstlisting}

\vspace{0.3em}
\textbf{Error Annotation.}
The first error blocks evaluation before the model is even built: the executor
sees the markdown fence token as Python syntax. The second error reaches
model construction but fails when a string shift name is used as a list index.
Both examples primarily motivate SFT: the model already has useful OR
knowledge, but Pass@1 behavior must obey code-only contracts and schema-safe
data-access templates.

\end{casepage}

\begin{casepage}{Case Study: Discrete Semantics Error in NL4OPT}

\begin{tabularx}{\linewidth}{@{}p{0.22\linewidth}X@{}}
\toprule
Benchmark & NL4OPT \\
Sample ID & \texttt{NL4OPT\_15} \\
Evaluation & Ground truth: 28250; Prediction: 28125 \\
Failure Type & Continuous relaxation of integer workforce variables \\
\bottomrule
\end{tabularx}

\vspace{0.5em}
\textbf{Original Problem}

{\small
A store employs senior citizens who earn \$500 per week and young adults who earn \$750 per week.
The store must keep the weekly wage bill below \$30000.
On any day, the store requires at least 50 workers, of whom at least 10 must be young adults.
To ensure the store runs smoothly, the number of young adults should be at least a third the number of senior citizens.
Formulate a LP to minimize the wage bill.
\par}

\vspace{0.5em}
\textbf{Model Output}

\begin{lstlisting}[style=pythoncase]
import gurobipy as gp
from gurobipy import GRB

m = gp.Model("MinimizeWageBill")

x = m.addVar(vtype=(*@\textbf{\textcolor{errRed}{GRB.CONTINUOUS}}@*), name="seniors")
y = m.addVar(vtype=(*@\textbf{\textcolor{errRed}{GRB.CONTINUOUS}}@*), name="young_adults")

m.setObjective(500 * x + 750 * y, sense=GRB.MINIMIZE)

m.addConstr(x + y >= 50, "total_workers")
m.addConstr(y >= 10, "min_young")
m.addConstr(3 * y >= x, "young_senior_ratio")
m.addConstr(500 * x + 750 * y <= 30000, "wage_budget")

m.optimize()

if m.status == GRB.OPTIMAL:
    print(f"Optimal solution: x = {x.X:.2f}, y = {y.X:.2f}")
    print(f"Minimum wage bill = ${m.ObjVal:.2f}")
else:
    print("No optimal solution found.")
\end{lstlisting}

\vspace{0.3em}
\textbf{Error Annotation.}
The problem asks for numbers of workers, but the model uses continuous variables.
This relaxes the integer workforce decision and yields a fractional optimum.

\end{casepage}

\begin{casepage}{Case Study: Loss/Yield-Driven Conservation Error}

\begin{tabularx}{\linewidth}{@{}p{0.22\linewidth}X@{}}
\toprule
Benchmark & B4O Feasible \\
Sample ID & \texttt{BENCH4OPT\_2} \\
Evaluation & Pass@1 prediction: 265.5768398591736; ground truth: 531.1536797183472. 5-shot and Pass@16 recover the reference value. \\
Failure Type & Evaporation loss mixed into both cost and flow conservation \\
\bottomrule
\end{tabularx}

\vspace{0.5em}
\textbf{Original Problem}

{\small
The task is a water-distribution network-flow problem. Each edge has a
capacity, transportation cost, and evaporation rate. The model must satisfy
node demand while minimizing the total cost induced by flow and evaporation.
\par}

\vspace{0.5em}
\textbf{Model Output}

\begin{lstlisting}[style=pythoncase]
flow = {}
for i, (source, target, capacity, cost, evaporation_rate) in enumerate(edges):
    flow[(source, target)] = model.addVar(lb=0.0, ub=capacity,
                                         name=f"flow_{source}_{target}")

obj = gp.LinExpr()
for (source, target), var in flow.items():
    ...
    obj += (*@\textbf{\textcolor{errRed}{\texttt{(cost + evaporation\_rate) * var}}}@*)
model.setObjective(obj, GRB.MINIMIZE)

for node_idx, node in enumerate(nodes):
    inflow = gp.LinExpr()
    outflow = gp.LinExpr()
    evaporation_loss = gp.LinExpr()
    ...
    model.addConstr(
        inflow - outflow - (*@\textbf{\textcolor{errRed}{\texttt{evaporation\_loss}}}@*) == demand[node_idx],
        name=f"flow_conservation_{node}"
    )
\end{lstlisting}

\vspace{0.3em}
\textbf{Error Annotation.}
The code is syntactically valid and reaches an optimal solution, but it
implements the wrong physical system. Evaporation is treated as an additive
unit cost and also deducted from conservation, blurring whether loss affects
delivered quantity, incurred cost, or both. This is a domain-pattern failure:
network flow with loss, yield, shrinkage, and evaporation needs stable
conservation templates rather than generic LP scaffolding.

\end{casepage}

\begin{casepage}{Case Study: Ratio, Unit, and Total-Quantity Coupling Error}

\begin{tabularx}{\linewidth}{@{}p{0.22\linewidth}X@{}}
\toprule
Benchmark & OptiBench \\
Sample ID & \texttt{OPTIBENCH\_6} \\
Evaluation & Pass@1 prediction: 10000.0; 5-shot prediction: 28000.0; Pass@16 prediction: 10000.0; ground truth: 18000.0 \\
Failure Type & Investment-return units and total savings are coupled incorrectly \\
\bottomrule
\end{tabularx}

\vspace{0.5em}
\textbf{Original Problem}

{\small
A logistics company allocates 50 trucks among three goods. Each good has an
initial per-truck delivery cost. Investment in route-optimization software can
reduce per-truck cost, and total investment must not exceed the total savings
generated by the investment. The objective is to minimize total delivery cost
while satisfying minimum truck counts for Goods X and Y.
\par}

\vspace{0.5em}
\textbf{Model Output}

\begin{lstlisting}[style=pythoncase]
x = m.addVars(goods, vtype=GRB.INTEGER, name="trucks")
I = m.addVars(goods, vtype=GRB.CONTINUOUS, name="investment")

# cost per truck = c_i - (r_i/100)*I_i
obj = gp.QuadExpr()
for i in goods:
    obj += (c[i] - (r[i]/100.0) * I[i]) * x[i]
m.setObjective(obj, GRB.MINIMIZE)

m.addConstr(gp.quicksum(I[i] for i in goods) <= B, name="budget")

# Savings constraint: total investment <= total cost savings
savings_lhs = gp.quicksum(I[i] for i in goods)
savings_rhs = gp.quicksum(
    (*@\textbf{\textcolor{errRed}{(r[i]/100.0) * I[i] * x[i]}}@*) for i in goods
)
m.addConstr(savings_lhs <= savings_rhs, name="savings")
m.params.NonConvex = 2
\end{lstlisting}

\vspace{0.3em}
\textbf{Error Annotation.}
The model recognizes that investment and truck allocation interact, but it does
not correctly reconcile dollars, hundreds-of-dollars, per-truck reduction, and
total savings. In 5-shot, the error changes shape: the model moves toward a
linear-looking savings template but can drop the scale variable. This case
therefore exposes a deeper semantic issue than API usage: ratio, ROI,
average-cost, and efficiency statements require explicit tracking of numerator,
denominator, unit, and total quantity.

\end{casepage}

\begin{casepage}{Case Study: Solver-Compatible Nonlinear Reformulation Failure}

\begin{tabularx}{\linewidth}{@{}p{0.22\linewidth}X@{}}
\toprule
Benchmark & OptiBench \\
Sample ID & \texttt{OPTIBENCH\_13} \\
Evaluation & Pass@1: optimization failed with ``Objective must be linear or quadratic''; 5-shot prediction: 0.125; Pass@16 prediction: 0.125; ground truth: 0.35355235026576626 \\
Failure Type & Fourth-order objective passed directly to Gurobi \\
\bottomrule
\end{tabularx}

\vspace{0.5em}
\textbf{Original Problem}

{\small
A person is trapped near a swamp boundary described by \(y \ge x^2\) with
\(x \ge 0\). The task is to find the closest feasible rescue point to
\((3/4,0)\), equivalently the shortest rope length to a point on or above the
parabolic boundary.
\par}

\vspace{0.5em}
\textbf{Model Output}

\begin{lstlisting}[style=pythoncase]
m = gp.Model("SwampRescue")
x = m.addVar(lb=0.0, vtype=GRB.CONTINUOUS, name="x")

# Objective: minimize (x - 0.75)^2 + (x^2)^2
obj = (x - 0.75)*(x - 0.75) + (*@\textbf{\textcolor{errRed}{x*x*x*x}}@*)
m.setObjective(obj, GRB.MINIMIZE)

m.optimize()
\end{lstlisting}

\vspace{0.3em}
\textbf{Error Annotation.}
The mathematical expression is meaningful, but it is not a linear or quadratic
objective. Later inference settings avoid the initial build error yet still
return the wrong distance, so the missing capability is not merely a Gurobi API
habit. The model needs solver-aware reformulation knowledge: when to introduce
auxiliary variables, when nonconvex quadratic or general constraints are
legitimate, and when analytic simplification is required before coding the
optimization model.

\end{casepage}